\chapter{Picard–Lefschetz 理论}\label{chapter:III}
\SourcePage{219}
本书所综合的第三项内容，是 Lefschetz 纤维化的辛几何。这是一个极其丰富的课题，我们的论述远不足以充分展现它。这里采用的定义基本与\bcite{131}相同。事实上，许多讨论都是对那篇论文的扩充或评述；这也意味着本书的处理有些偏于一面。有兴趣获得更全面认识的读者，可以从\bcite{40}、\bcite{64}、\bcite{11}、\bcite{17}入手。\znote{211}{Picard–Lefschetz 理论把全纯函数奇点的研究与其定义域复流形的拓扑联系起来，可视为复几何中与 Morse 理论平行的一套方法。其出发点是 Picard–Lefschetz 公式：临界点附近局部形如 $z_1^2+\cdots+z_n^2$，绕过临界值的环路诱导的单值变换，在消失圈上表现为一次 Dehn 扭转（相差一个符号）。由此，纤维化的单值变换可以完全用消失圈的相交数描述，这正是本章把 Lefschetz 纤维化问题化为组合数据的基础。[来源：nLab "Picard-Lefschetz theory"，https:/\allowbreak /\allowbreak ncatlab.org/\allowbreak nlab/\allowbreak show/\allowbreak Picard-Lefschetz+\allowbreak theory（2026-11-15 访问）。]}

前两节回顾一些初等方面。与\bcite{131}相比，主要增加的是 Donaldson 提出的匹配圈概念，另见\bcite{16}。粗略地说，若某个 Lefschetz 纤维化的一般纤维本身又是第二个 Lefschetz 纤维化的总空间，就会出现匹配圈。我们把这种结构称为双纤维化；在最低的非平凡维数中，它是\bcite{12}、\bcite{13}所考察的辛分歧覆盖的一种较简单的版本。\znote{211}{“一般纤维本身又是总空间”意味着有链 $E\to S\to W$，其中 $E\to S$ 与 $S\to W$ 都是带奇点的纤维化。此时固定 $w\in W$，$E_w\to S_w$ 又给出一个纤维化，而 $S_w$ 是曲面、$E_w$ 比它高两维。匹配圈就是在这个结构里构造的 Lagrangian 球面：在 $S_w$ 中取一条连接两个临界值的路径，把两端纤维中的消失圈沿路径平行移动并粘合。路径的选取与粘合方式不同，所得球面可能不同，这正是 Donaldson 的构造所利用的自由度。[来源：原书 §(15e) 双纤维化的定义，印刷页 218；匹配圈见 \bcite{16}。]}

随后，考虑以 Lefschetz 纤维化的全纯截面为基础的 Floer TQFT 扩展，并借此导出一个基本比较定理，联系第\sr{sec:5}{5 节}的代数理论与第\sr{sec:16}{16 节}的辛几何。这个定理在对象层次将扭转函子 $T$ 的作用与辛 Dehn 扭转 $\tau$ 的作用认同；事实上，它是\bcite{131}中长正合列的另一种表述。我们不重述该正合列的完整证明，但会较详细地解释，如何把结果变成这里所需的形式。

给定以圆盘为底空间的 Lefschetz 纤维化 $\pi:E\to D$，可以取沿某个纤维 $E_z$ 分歧的二重覆盖；代数几何中称之为“底扩张”。实际上，改变 $z$ 就得到一族这样的覆盖，这是双纤维化的一个特别简单的例子。二重覆盖显然有 $\Z/2$ 作用，可对它应用第\sr{sec:6}{6 节}和第\sr{sec:14}{14 节}的等变技术。我们将用这一想法定义 $\Ai$ 范畴 $\F(\pi)$，称为 \emph{Lefschetz 纤维化的 Fukaya 范畴}，它是纤维 $E_z$ 与总空间 $E$ 的 Fukaya 范畴之间的中介。诚然，借二重覆盖技巧给出的定义并非最自然，但它使大量已有工具得以重新使用。

接下来考察经典代数几何意义下的 Lefschetz 束，并证明引言中陈述的结果。最后一节讨论来自代数几何中一个稍有不同领域的一系列例子，即 $(A_m)$ 型简单奇点的 Milnor 纤维。它们首先具有历史意义，因为这是最早有人为导出 Fukaya 范畴提出显式代数模型的情形之一\bcite{85}；后来发现，那些模型只在实维数 $\geq4$ 时正确。到这里，讨论相较全书其余部分不再那么自足，但希望能借此展示具体计算中使用的几何与代数技巧。关于同一套一般工具的更复杂
\SourcePage{220}
应用，请读者参看文献，例如\bcite{14}、\bcite{15}、\bcite{130}、\bcite{59}。
